%% kapa: exam=1 round=26 year=2015 period=2 form=A subject=appraisal_law subject_ko=부동산관계법규 seq=16 number=16
%% source: Q-Net 감정평가사 기출문제 (cst00309) artlSeq=5138001 file=extracted/1008/5138001/제26회_감정평가사 1차_기출/제26회_감정평가사 1차_기출/2교시_감평_A형.pdf
%% conversion: kordoc 4.24.51 -> tools/gen1.py (auto); review=proofread-sonnet; flags=none
\begin{question}{16}
국토의 계획 및 이용에 관한 법령상 국가계획, 광역도시계획 및 도시ㆍ군계획의 관계에 관한 설명으로 옳지 않은 것은?
\begin{choices}
  \item 광역도시계획의 내용이 국가계획의 내용과 다를 때에는 국가계획의 내용이 우선한다.
  \item 도시ㆍ군기본계획의 내용이 광역도시계획의 내용과 다를 때에는 광역도시계획의 내용이 우선한다.
  \item 도시ㆍ군계획은 특별시ㆍ광역시ㆍ특별자치시ㆍ특별자치도ㆍ시 또는 군의 관할 구역에서 수립되는 다른 법률에 따른 토지의 이용ㆍ개발 및 보전에 관한 계획의 기본이 된다.
  \item 도시ㆍ군계획은 국가계획에 부합되어야 하며, 도시ㆍ군계획의 내용이 국가계획의 내용과 다를 때에는 국가계획의 내용이 우선한다.
  \item 특별시장이 관할 구역에 대하여 다른 법률에 따른 환경ㆍ교통ㆍ수도ㆍ하수도ㆍ주택에 관하여 수립하는 부문별 계획은 도시ㆍ군기본계획의 내용에 부합되지 않아도 된다.
\end{choices}
\end{question}
\answer{5}
